\subsection{مبانی نظری و فرمول‌بندی ریاضی}
در گام نخست، ماتریس کوواریانس تجربی فعال‌سازی‌های ورودی $\mathbf{C} \in \mathbb{R}^{d_{\text{in}} \times d_{\text{in}}}$ با عبور دسته‌های کالیبراسیون از مدل تخمین زده می‌شود:
\begin{equation}
\mathbf{C} = \frac{1}{N} \sum_{k=1}^N \mathbf{X}_k^\top \mathbf{X}_k
\end{equation}
جهت رفع تکینگی و تضمین قطعی معین مثبت بودن، ماتریس با ضریب میراگر $\lambda_{\text{damp}}$ و میانگین درایه‌های قطری منظم‌سازی می‌شود:
\begin{equation}
\mathbf{C}_{\text{reg}} = \mathbf{C} + \max\left(\lambda_{\text{damp}} \cdot \text{mean}(\text{diag}(\mathbf{C})), 10^{-6}\right) \mathbf{I}
\end{equation}
سپس تجزیه مقادیر ویژه متقارن بر روی $\mathbf{C}_{\text{reg}}$ اجرا می‌شود:
\begin{equation}
\mathbf{C}_{\text{reg}} = \mathbf{Q} \mathbf{\Lambda} \mathbf{Q}^\top, \quad \mathbf{\Lambda} = \text{diag}(\lambda_1, \dots, \lambda_{d_{\text{in}}})
\end{equation}
با برش مقادیر ویژه به آستانه عددی $\lambda_i \ge 10^{-6}$، جذر ماتریس و معکوس جذر آن به صورت متقارن محاسبه می‌گردند:
\begin{equation}
\mathbf{C}^{1/2} = \mathbf{Q} \mathbf{\Lambda}^{1/2} \mathbf{Q}^\top, \quad \mathbf{C}^{-1/2} = \mathbf{Q} \mathbf{\Lambda}^{-1/2} \mathbf{Q}^\top
\end{equation}
ماتریس وزن اولیه لایه $\mathbf{W} \in \mathbb{R}^{d_{\text{out}} \times d_{\text{in}}}$ با ضرب در جذر کوواریانس سفیدسازی می‌شود:
\begin{equation}
\mathbf{M} = \mathbf{W} \mathbf{C}^{1/2} \in \mathbb{R}^{d_{\text{out}} \times d_{\text{in}}}
\end{equation}
سپس تجزیه مقادیر منفرد تقلیل‌یافته با رتبه $r$ بر روی $\mathbf{M}$ اعمال می‌شود:
\begin{equation}
\mathbf{M} \approx \mathbf{U}_r \mathbf{\Sigma}_r \mathbf{V}_r^\top
\end{equation}
ماتریس بازسازی‌شده سفیدشده $\mathbf{M}_r = \mathbf{U}_r \mathbf{\Sigma}_r \mathbf{V}_r^\top$ از طریق معکوس جذر کوواریانس دی‌سفیدسازی (\lr{De-whitening}) می‌گردد:
\begin{equation}
\widehat{\mathbf{W}} = \mathbf{M}_r \mathbf{C}^{-1/2}
\end{equation}
خطای نسبی بازسازی فروبنیوس ارزیابی می‌شود:
\begin{equation}
\mathcal{E} = \frac{\|\mathbf{W} - \widehat{\mathbf{W}}\|_F}{\|\mathbf{W}\|_F}
\end{equation}
اگر $\mathcal{E} \le \tau_{\text{max}}$ باشد، وزن تقریب‌زده شده جایگزین می‌گردد؛ در غیر این صورت، ماژول ابتدا به تجزیه مقادیر منفرد استاندارد (\lr{Standard SVD Fallback}) و در صورت عدم برآورده شدن کیفیت، به وزن‌های اولیه رجوع می‌کند.

که در آن $\sigma_i(\mathbf{W})$ بیانگر $i$-امین مقدار منفرد مرتب‌شده ماتریس به صورت نزولی است. نسبت حفظ انرژی طیفی $\eta_r = \frac{\sum_{i=1}^r \sigma_i^2}{\sum_{i=1}^{d} \sigma_i^2}$ تضمین می‌کند که با انتخاب آستانه $\tau \ge 0.90$، بیش از نود درصد از پراکندگی واریانس داده‌ها حفظ شده و افت سیگنال در انتشار گرادیان رو به عقب به حداقل ممکن کاهش می‌یابد.

\paragraph{تحلیل مجانبی و پیچیدگی اشتراک حافظه سخت‌افزاری}
از منظر پیچیدگی محاسباتی و انتقال داده‌ها در سلسومراتب حافظه (حافظه سراسری \lr{HBM} به حافظه محلی \lr{SRAM})، عملیات ضرب ماتریسی در این مکانیزم نیازمند تخصیص بافر موقت با ابعاد فاکتورگیری‌شده است. برای پردازش دنباله‌ای به طول $N$ با بعد مخفی $D$، بار محاسباتی هسته از فرم درجه دو به رابطه زیر تقلیل می‌یابد:
\begin{equation}
\label{eq:complexity_gfws_d_compressor}
\mathcal{C}_{\text{gfws\_d\_compressor}} = \mathcal{O}\left( B \cdot N \cdot D \cdot r + B \cdot H \cdot N \cdot d_k \right) \ll \mathcal{O}\left( B \cdot H \cdot N^2 \cdot d_k \right)
\end{equation}
این بهبود مرتبه پیچیدگی، با کاهش چشمگیر حجم تبادل تانسورها در گذرگاه‌های ارتباطی کارت گرافیک، مانع از اشباع پهنای باند حافظه (پدیده \lr{Memory-Bound}) در فرآیند استنتاج طولانی‌مدت می‌گردد.
